\subsection{Criterios de selección}

Con el fin de elegir un simulador adecuado para integrar y evaluar compresión \emph{sin pérdida de precisión} sobre el vector de estado, se definieron los criterios de la tabla \ref{tab:criterios_sim}. Estos criterios operacionalizan los lineamientos metodológicos del diseño (modificabilidad, integración de estructuras, optimización/paralelismo e instrumentación), y además incorporan requisitos prácticos de ingeniería para reducir riesgo técnico durante la implementación.

\begin{table}[H]
    \centering
    \small
    \setlength{\tabcolsep}{3pt}
    \vspace{-2mm}
    \caption{Criterios usados para la selección del simulador cuántico} \label{tab:criterios_sim}
    \begin{tabular}{>{\raggedright\arraybackslash}p{0.7cm}>{\raggedright\arraybackslash}p{4.0cm}>{\raggedright\arraybackslash}p{0.56\linewidth}>{\centering\arraybackslash}p{0.9cm}}
        \hline
        \multicolumn{1}{l}{} &
        \multicolumn{1}{c}{\textbf{Criterio}} &
        \multicolumn{1}{c}{\textbf{Explicación}} &
        \multicolumn{1}{c}{\textbf{Peso}} \\ \hline

        C1 & Vector de estado apto para compresión sin pérdida &
        El simulador debe representar el vector de estado como arreglos contiguos de valores de punto flotante en \textbf{doble precisión} (amplitudes complejas), de manera que la compresión sin pérdida se pueda aplicar directamente sobre dicha estructura sin cambios semánticos. &
        0.10 \\

        C2 & Facilidad de integración de estructuras o arreglos personalizados &
        El código debe permitir reemplazar o envolver el almacenamiento del vector de amplitudes (p.\ ej.\ insertar un contenedor comprimido) con impacto controlado en el resto del simulador. Este criterio prioriza código modular y rutas críticas identificables. &
        0.25 \\

        C3 & Soporte y afinidad con optimización de memoria &
        Se valora que el simulador esté concebido para experimentar con estrategias de memoria (p.\ ej.\ compresión, poda/pruning, particionamiento o variantes del almacenamiento), reduciendo esfuerzo para insertar y validar la estrategia propuesta. &
        0.20 \\

        C4 & Complejidad de construcción y dependencias &
        Un proceso de compilación simple y pocas dependencias facilitan reproducibilidad y portabilidad de experimentos. Se penalizan toolchains complejas cuando incrementan el esfuerzo de integración e instrumentación. &
        0.10 \\

        C5 & Capacidad de instrumentación (tiempo y memoria) &
        El simulador debe permitir medir de forma fina (por secciones) el consumo de memoria y el tiempo de ejecución, idealmente en puntos cercanos al acceso al vector de estado. &
        0.15 \\

        C6 & Cobertura funcional mínima para casos de prueba &
        Debe ser posible ejecutar circuitos estándar de evaluación (p.\ ej.\ QFT y/o Grover) con un conjunto mínimo de compuertas, para poder comparar desempeño con y sin compresión. &
        0.10 \\

        C7 & Paralelismo relevante (CPU, OpenMP, MPI) &
        Se considera valioso poder evaluar la interacción compresión--paralelismo. Se prioriza soporte a OpenMP/MPI cuando el objetivo incluye escalamiento o comparación bajo concurrencia. &
        0.10 \\ \hline
    \end{tabular}
    \vspace{-3mm}
\end{table}

Los pesos fueron definidos para privilegiar la \textbf{integración del almacenamiento del vector de estado} (C2) y la \textbf{afinidad con experimentación en memoria} (C3), dado que son los factores que más reducen el riesgo técnico de insertar compresión sin pérdida sin distorsionar el comportamiento del simulador.
